% Part: first-order-logic
% Chapter: introduction
% Section: soundness-completeness

\documentclass[../../../include/open-logic-section]{subfiles}

\begin{document}

\olfileid{fol}{int}{scp}

\olsection{سموالے او بشپړتيا}

د لومړۍ درجې منطق دپاره به د !!{derivation} نظامونه هم ورپېژنو۔
منطقپوهانو ډېر د !!{derivation} نظامونه جوړ کړي دي، خو هغوئ ټول د
!!{sentence}s تر منځ هماغه د !!{derivability} اړيکه تعريفوي۔ وايو
چې~$\Gamma$،~$!A$ !!{derive}s کوي، $\Gamma \Proves !A$، که د يوه
ټاکلي، په دقيق ډول تعريف شوي ډول !!a{derivation} موجود وي۔
!!^{derivation}s تل د سمبولونو متناهي ترتيبونه دي---ښايي د
!!{sentence}s لړ وي، يا کوم لا پېچلے جوړښت۔ د !!{derivation} د
نظامونو موخه داسې وسيله برابرول دي چې وټاکي ايا کوم سټ~$\Gamma$
!!a{sentence} لازموي که نۀ۔ د دې موخې د پوره کولو دپاره بايد
$\Gamma \Entails !A$ هله او يوازې هله وي چې $\Gamma \Proves !A$
وي۔

که $\Gamma \Proves !A$ وي خو $\Gamma \Entails !A$ نۀ وي، زموږ د
!!{derivation} نظام به ډېر پياوړے وي او له حده زيات څه به ثابتوي۔
هغه خاصيت چې له~$\Gamma \Proves !A$ څخه $\Gamma \Entails !A$
لازمېږي، \emph{سموالے} بلل کېږي، او د !!{derivation} د هر ښه نظام
لږ تر لږه غوښتنه ده۔ له بلې خوا، که $\Gamma \Entails !A$ وي خو
$\Gamma \Proves !A$ نۀ وي، نو زموږ د !!{derivation} نظام ډېر کمزورے
دے او کافي څه نۀ ثابتوي۔ هغه خاصيت چې له~$\Gamma \Entails !A$ څخه
$\Gamma \Proves !A$ لازمېږي، \emph{بشپړتيا} بلل کېږي۔ سموالے
عموماً نسبتاً اسانه ثابتېږي (د !!{derivation}s پر جوړښت د استقرا له
لارې، ځکه هغوئ په استقرا تعريف شوي دي)۔ بشپړتيا ثابتول ګران دي۔

سموالے او بشپړتيا څو مهمې پايلې لري۔ که د !!{sentence}s يو
سټ~$\Gamma$ کوم تناقض (لکه~$!A \land \lnot !A$) !!{derive}s کړي،
\emph{ناسازګار} بلل کېږي۔ ناسازګار~$\Gamma$ هېڅ مدل نۀ شي لرلے؛ د
صدق وړ نۀ دے۔ له بشپړتيا څخه ئې عکس لازمېږي: هر هغه~$\Gamma$ چې
ناسازګار نۀ وي---يا، لکه موږ چې به وايو، \emph{سازګار} وي---يو مدل
لري۔ په حقيقت کښې دا له بشپړتيا سره هم‌ارزښته ده، او موږ به د
بشپړتيا همدا بڼه په اصل کښې ثابتوو۔ دا ژوره او ښايي حيرانوونکې
پايله ده: يوازې دا چې له~$\Gamma$ څخه~$!A \land \lnot !A$ نۀ شئ
ثابتولے، تضمينوي چې داسې !!a{structure} شته چې کټ مټ د~$\Gamma$ د
بيان په شان ده۔ نو بشپړتيا دې پوښتنې ته ځواب ورکوي: د
!!{sentence}s کوم سټونه مدلونه لري؟ ځواب: ټول او يوازې سازګار
سټونه۔

د سموالي او بشپړتيا قضيې دوه مهمې پايلې لري: د متناهي شاهد قضيه او
د لوېنهايم--سکولم قضيه۔ دا د مدلونو په تيورۍ کښې مهمې پايلې دي
او د ډېرو په زړه پورې پايلو د ثبوت دپاره کارېدے شي۔ دوه مو لا يادې
کړې دي: د لومړۍ درجې منطق دا نۀ شي بيانولے چې د !!a{structure}
!!{domain} متناهي يا !!{nonenumerable} ده۔

په تاريخي ډول د دې ټولو کارونو سم معلومول ډېر وخت ونيو---د لومړۍ
درجې منطق نحو او معنٰی پوهنه څنګه تعريف شي، د !!{derivation} ښه
نظامونه څنګه تعريف شي، د هغوئ سموالے او بشپړتيا څنګه ثابته شي، او
د لومړۍ درجې ژبو د بيانېدونکو او نۀ بيانېدونکو څيزونو توپير څنګه
روښانه شي۔ اوس پوهېږو چې دا هر څه څنګه وکړو، خو په ټولو جزئياتو
تېرېدل لا هم مغشوشوونکي او ستومانوونکي کېدے شي۔ بيا هم مهم دي، ځکه
دلته د لومړۍ درجې منطق د صوري ژبې دپاره جوړې شوې طريقې په منطق،
کمپيوټر ساينس او ژبپوهنه کښې هر ځاے پلي کېږي۔ نو د جزئياتو زيار په
اوږده موده کښې ګټه کوي۔

\end{document}
